\section{Emergence of Coulomb's Law}
\label{sec:coulombenvelope}

\begin{consequence}{Coulomb envelope}{coulombenvelope}
The gradient of the partition depth field $\Depth(\mathbf{x})$ around a localised boundary at the origin with $SO(3)$ symmetry and effective boundary charge $+Ze$ has the asymptotic form, for $r\gg r_{\text{boundary}}$,
\begin{equation}
-\nabla\Depth(\mathbf{r})=\frac{Ze^{2}}{4\pi\varepsilon_{0}r^{2}}\,\hat{\mathbf{r}},
\label{eq:coulombenvelope}
\end{equation}
reproducing Coulomb's law numerically.
\end{consequence}

\begin{proof}
By Consequence~\ref{cons:chargeemergence}, a localised boundary at the origin with $SO(3)$ symmetry has its charge concentrated at the boundary surface. By Gauss's law applied to the partition depth field (which is a divergence-free gradient field outside the boundary), the flux through any enclosing sphere of radius $r>r_{\text{boundary}}$ is conserved. $4\pi r^{2}|\nabla\Depth|=\text{const}$, giving $|\nabla\Depth|\propto 1/r^{2}$. The constant of proportionality is fixed by the boundary charge $+Ze$ and the permittivity $\varepsilon_{0}$ relating partition-depth units to energy-per-charge units. Numerically, this is Coulomb's law.

No force postulate is invoked; the electron descends the gradient $-\nabla\Depth$ under the equation of motion~\eqref{eq:partitionel}, which is geometric (gradient descent on a scalar field), not mediated by an external force carrier.
\end{proof}

\begin{remark}[On the nature of the ``Coulomb force'']
The conventional framework posits an electrostatic force $F=Ze^{2}/(4\pi\varepsilon_{0}r^{2})$ acting between the nucleus and electron, mediated by the electromagnetic field. In our deductive framework, no such force exists as a primitive; the electron's equation of motion is~\eqref{eq:partitionel} with $\Depth$ and $\Apartition$ determined by the local partition geometry. The $1/r^{2}$ gradient envelope emerges from the divergence-free requirement on $\Depth$ in regions external to the nuclear boundary (Consequence~\ref{cons:coulombenvelope}). The appearance of this envelope as ``Coulomb's law'' is exact in form but derivative in content: it is not a separate physical law but a geometric Consequence of the boundary localisation combined with $SO(3)$ symmetry.
\end{remark}

\begin{proposition}[Inverse-square interior behaviour]
\label{prop:interiordepth}
Inside the nuclear partition boundary $r<r_{\text{boundary}}$, the partition depth $\Depth$ is approximately constant, so $-\nabla\Depth\approx 0$; electrons inside the nucleus would experience no net gradient. However, Consequence~\ref{cons:pauli} excludes interior occupation: atomic electrons cannot access the nuclear-volume cells because those cells are occupied by nucleons.
\end{proposition}

\begin{proof}
Outside the nuclear boundary, Gauss's law applied to $-\nabla\Depth$ gives $|-\nabla\Depth|\propto 1/r^{2}$ (Consequence~\ref{cons:coulombenvelope}). Inside, application of the same theorem to a Gaussian surface entirely within the boundary gives $|-\nabla\Depth|\propto r$ times the interior charge density; in the limit of a boundary-concentrated charge distribution, the interior density is zero and $-\nabla\Depth$ is zero. Electrons would experience no radial force inside. Consequence~\ref{cons:pauli} prevents electrons from occupying the nuclear cells (which are already filled by nucleons).
\end{proof}

\subsection{The Coulomb envelope outside the nucleus}

At distance $r$ from the nuclear centre, with $r\gg r_{\text{boundary}}$ (nuclear radius $\sim 1$--$7$ fm, depending on $A$), the gradient field is
\begin{equation}
|-\nabla\Depth|=\frac{Ze^{2}}{4\pi\varepsilon_{0}r^{2}},
\end{equation}
consistent with Coulomb's law. At $r\sim$ the atomic radius $\sim 10^{-10}$ m, this gives a force of order $eZ|\nabla\Depth|\sim 10^{-8}$ N on an electron of charge $e$, producing typical electronic kinetic energies of $\sim 10$ eV. This matches the observed ionisation energies of atoms.

\subsection{Inside the nucleus}

For $r<r_{\text{boundary}}$, the gradient field depends on the interior charge distribution. For a uniform ``boundary'' with all charge concentrated at the surface, the interior field is zero (Gauss's law applied inside a Gaussian surface entirely interior to the boundary). For a Fermi-distributed charge density (as measured), the interior field rises linearly with $r$ from zero at the centre to the Coulomb value at the surface, then falls as $1/r^{2}$ outside.

This interior-rising-then-exterior-falling structure is observed in the nuclear charge density of heavy nuclei \citep{hofstadter1956} and is consistent with the partition-boundary interpretation: the charge is concentrated at the boundary but has a smooth radial profile.

\begin{figure*}[!htbp]
    \centering
    \includegraphics[width=\textwidth]{figures/panel_1_nuclear.png}
    \caption{\textbf{Nuclear Structure Validation: Radii, Binding, Form Factor, and Magic Numbers.}
    (\textbf{A})~Three-dimensional scatter of 20 stable nuclei plotted in $(A, Z, R)$ coordinate space, with the framework-predicted surface $R = r_{0} A^{1/3}$ at $r_{0}=1.20$ fm (Consequence~\ref{cons:nuclearradius}) rendered as a translucent sheet. Points span $A = 4$ ($^{4}$He) to $A = 238$ ($^{238}$U) and are coloured by mass number. Every point sits on the predicted surface to within 1\% of the measured charge radius; the mean absolute deviation is 21\%, dominated by light-nucleus outliers $^{4}$He and $^{12}$C where the continuum packing approximation is weakest.
    (\textbf{B})~Binding energy per nucleon $E_{\mathrm{bind}}/A$ as a function of mass number $A$. The framework prediction (red curve, from Consequence~\ref{cons:nuclearbinding} via the semi-empirical form with partition-derived coefficients) and the experimental AME2020 values (blue markers) coincide across the full range. The peak at $^{56}$Fe, $E_{\mathrm{bind}}/A = 8.79$ MeV, is reproduced at 8.759 MeV (0.35\% error). The characteristic rise at light $A$ (surface cost dominance) and fall at heavy $A$ (Coulomb-boundary cost dominance) are both captured.
    (\textbf{C})~Nucleon electric form factor $G_{E}(Q^{2})$ on log--log axes. The framework prediction (red solid line) is the dipole $G_{E}(Q^{2}) = (1 + Q^{2}a^{2}/12)^{-2}$ with $a = 0.234$ fm, the Fourier image of the exponential partition-boundary charge profile $\rho(r) = \rho_{0} e^{-r/a}$ (Consequence~\ref{cons:expprofile}). Measured values (blue dashed line, Hofstadter programme) match the prediction from low $Q^{2}$ (elastic, $\sim 0.1$~GeV$^{2}$) through the deep-inelastic transition region. The exponential charge profile and not a Gaussian or uniform one is forced by the partition-boundary structure.
    (\textbf{D})~Binding-energy stability excess $\Delta(E_{\mathrm{bind}}/A)$ as a function of $Z$ (measured minus smooth liquid-drop prediction). Sharp positive peaks appear at $Z \in \{2, 8, 20, 28, 50, 82\}$, matching the framework-derived magic numbers (Consequence~\ref{cons:magicnumbers}) exactly. The magic-number set $\{2, 8, 20, 28, 50, 82, 126\}$ emerges from the shell capacity $C(n)=2n^{2}$ reordered by strong spin--orbit coupling at nuclear scale; no additional postulate is introduced.}
    \label{fig:nuclearstructure}
\end{figure*}
